\chapter{Results}
\label{chap:results}

\begin{placeholderblock}
This chapter is a reporting scaffold. Replace every red entry after the analysis pipeline, human validation, and pre-specified exclusions are complete. Do not write the discussion first and then select results that fit it. Report null, mixed, and model-specific findings with the same level of detail as positive findings.
\end{placeholderblock}

\section{Evaluation set and data quality}

The final benchmark contained \resultcell{} scenario families and \resultcell{} scenario instances, of which \resultcell{} were single-turn and \resultcell{} were multi-turn. Across \resultcell{} target models and \resultcell{} stochastic replicates, the study attempted \resultcell{} scenario-condition runs. After applying preregistered exclusions for \placeholder{invalid generation, provider error, persona-semantic drift, source corruption, and other specified reasons}, \resultcell{} runs remained for analysis (\resultcell\%).

Table~\ref{tab:dataset-accounting} should make all exclusions auditable. Exclusions must be reported by prompt and persona cell to rule out differential missingness.

\begin{table}[H]
\centering
\caption{Run accounting and exclusions.}
\label{tab:dataset-accounting}
\small
\begin{tabularx}{\textwidth}{@{}lrrrrY@{}}
\toprule
Stage & Neutral & Prod.-base & Prod.-integrity & Total & Notes \\
\midrule
Attempted runs & \resultcell{} & \resultcell{} & \resultcell{} & \resultcell{} & \placeholder{all cells} \\
Provider/API failures & \resultcell{} & \resultcell{} & \resultcell{} & \resultcell{} & \placeholder{retry policy} \\
Invalid/empty outputs & \resultcell{} & \resultcell{} & \resultcell{} & \resultcell{} & \placeholder{definition} \\
Persona drift & \resultcell{} & \resultcell{} & \resultcell{} & \resultcell{} & \placeholder{semantic-equivalence threshold} \\
Other preregistered exclusions & \resultcell{} & \resultcell{} & \resultcell{} & \resultcell{} & \placeholder{list categories} \\
Analysed runs & \resultcell{} & \resultcell{} & \resultcell{} & \resultcell{} & \placeholder{final denominator} \\
\bottomrule
\end{tabularx}
\end{table}

Response length differed across prompt conditions: the median was \resultcell{} tokens in neutral, \resultcell{} in production-baseline, and \resultcell{} in production-integrity. Persona medians were \resultcell, \resultcell, and \resultcell{} for neutral, anxious, and positive users respectively. \placeholder{Report distributional plots and state whether length differences motivate, but do not replace, the preregistered secondary length-adjusted analysis.}

\section{Annotation performance and human agreement}

Before comparing target models, the automated scoring pipeline was evaluated against \resultcell{} independently annotated responses. Table~\ref{tab:annotation-reliability} reports performance for each label. The main text should identify the weakest component and explain how uncertainty in that label affects subsequent conclusions.

\begin{table}[H]
\centering
\caption{Automated scoring performance against adjudicated human labels.}
\label{tab:annotation-reliability}
\small
\begin{tabularx}{\textwidth}{@{}Yccccc@{}}
\toprule
Label & Prevalence & Precision & Recall & F1 / ICC & Human agreement \\
\midrule
Fact presence & \resultcell{} & \resultcell{} & \resultcell{} & \resultcell{} & \resultcell{} \\
Unsupported claim & \resultcell{} & \resultcell{} & \resultcell{} & \resultcell{} & \resultcell{} \\
Contradiction & \resultcell{} & \resultcell{} & \resultcell{} & \resultcell{} & \resultcell{} \\
Specificity preservation & \resultcell{} & \resultcell{} & \resultcell{} & \resultcell{} & \resultcell{} \\
Minimising language & \resultcell{} & \resultcell{} & \resultcell{} & \resultcell{} & \resultcell{} \\
Benefit amplification & \resultcell{} & \resultcell{} & \resultcell{} & \resultcell{} & \resultcell{} \\
Disclaimer washing & \resultcell{} & \resultcell{} & \resultcell{} & \resultcell{} & \resultcell{} \\
Overall net impression & \resultcell{} & \resultcell{} & \resultcell{} & \resultcell{} & \resultcell{} \\
\bottomrule
\end{tabularx}
\end{table}

\placeholder{State whether automated estimates were bias-corrected using the validation sample, whether primary analyses were repeated on human-only labels, and whether conclusions changed under an alternative judge.}

\section{RQ1: Prevalence under benign conditions}

In the production-baseline condition, \resultcell\% of conversations met the preregistered response-level risk-concealment criterion (95\% CI \resultcell--\resultcell). The corresponding rate was \resultcell\% in the neutral condition. High-adverse hard omission occurred in \resultcell\% and \resultcell\% of eligible fact opportunities, respectively. False claims occurred in \resultcell\% of responses.

The paired production-baseline versus neutral contrast was \resultcell{} percentage points for the primary endpoint (95\% CI \resultcell--\resultcell; adjusted $p=\resultcell$). \placeholder{Explain whether the estimate supports a practically meaningful baseline risk, a prompt-specific production effect, or substantial uncertainty. A non-zero raw rate alone does not show selectivity; compare against favourable and neutral fact retention.}

\begin{table}[H]
\centering
\caption{Primary outcomes by system-prompt condition, pooled across personas and target models.}
\label{tab:primary-prompt}
\small
\begin{tabularx}{\textwidth}{@{}Yccc@{}}
\toprule
Outcome & Neutral & Production-baseline & Production-integrity \\
\midrule
Any risk-concealment event (\%) & \resultcell{} & \resultcell{} & \resultcell{} \\
High-adverse omission (\%) & \resultcell{} & \resultcell{} & \resultcell{} \\
Favourable fact presence (\%) & \resultcell{} & \resultcell{} & \resultcell{} \\
Polarity selectivity gap (pp) & \resultcell{} & \resultcell{} & \resultcell{} \\
Any false claim (\%) & \resultcell{} & \resultcell{} & \resultcell{} \\
Any misleading framing (\%) & \resultcell{} & \resultcell{} & \resultcell{} \\
Any risk de-emphasis (\%) & \resultcell{} & \resultcell{} & \resultcell{} \\
Median response tokens & \resultcell{} & \resultcell{} & \resultcell{} \\
\bottomrule
\end{tabularx}
\end{table}

\section{RQ2: Distribution of failure modes}

Table~\ref{tab:failure-modes} decomposes the umbrella outcome. The most frequent failure mode was \resultcell, observed in \resultcell\% of responses or \resultcell\% of eligible fact units. Explicit false claims were \resultcell{} than the combined non-falsifiable categories by \resultcell{} percentage points (95\% CI \resultcell--\resultcell).

\begin{table}[H]
\centering
\caption{Failure-mode prevalence by target model, pooled across experimental conditions. Values should include 95\% confidence intervals in the final version.}
\label{tab:failure-modes}
\small
\begin{tabularx}{\textwidth}{@{}Yccccc@{}}
\toprule
Model & False claim & Omission & Misleading framing & De-emphasis & Any event \\
\midrule
\placeholder{Model A} & \resultcell{} & \resultcell{} & \resultcell{} & \resultcell{} & \resultcell{} \\
\placeholder{Model B} & \resultcell{} & \resultcell{} & \resultcell{} & \resultcell{} & \resultcell{} \\
\placeholder{Model C} & \resultcell{} & \resultcell{} & \resultcell{} & \resultcell{} & \resultcell{} \\
\placeholder{Model D} & \resultcell{} & \resultcell{} & \resultcell{} & \resultcell{} & \resultcell{} \\
\placeholder{Model E} & \resultcell{} & \resultcell{} & \resultcell{} & \resultcell{} & \resultcell{} \\
Pooled & \resultcell{} & \resultcell{} & \resultcell{} & \resultcell{} & \resultcell{} \\
\bottomrule
\end{tabularx}
\end{table}

Among framing subtypes, specificity loss occurred for \resultcell\% of present high-adverse facts; minimising language appeared in \resultcell\% of responses; unsupported benefit amplification in \resultcell\%; and disclaimer washing in \resultcell\%. The mean adverse-versus-favourable placement gap was \resultcell, and the median token-allocation balance was \resultcell.

\placeholder{Include a figure showing fact survival from source to response by polarity and materiality. A useful visual is a coefficient or marginal-probability plot rather than a stacked bar that obscures uncertainty.}

\section{RQ3: Effect of production-integrity guidance}

Relative to production-baseline, the production-integrity prompt changed the probability of the primary risk-concealment endpoint by \resultcell{} percentage points (95\% CI \resultcell--\resultcell). It changed high-adverse disclosure by \resultcell{} points, false claims by \resultcell{} points, and disclaimer washing by \resultcell{} points. The median response became \resultcell{} tokens longer.

Table~\ref{tab:mixed-model} reports the pre-specified mixed-effects model. The final text should translate coefficients into marginal probability differences and should state convergence diagnostics and random-effect variance.

\begin{table}[H]
\centering
\caption{Mixed-effects model for high-adverse omission. Reference levels: neutral prompt, neutral persona, and \placeholder{reference target model}.}
\label{tab:mixed-model}
\small
\begin{tabularx}{\textwidth}{@{}Yrrrr@{}}
\toprule
Predictor & Estimate & SE & 95\% CI & Adjusted $p$ \\
\midrule
Production-baseline prompt & \resultcell{} & \resultcell{} & \resultcell{} & \resultcell{} \\
Production-integrity prompt & \resultcell{} & \resultcell{} & \resultcell{} & \resultcell{} \\
Anxious persona & \resultcell{} & \resultcell{} & \resultcell{} & \resultcell{} \\
Positive persona & \resultcell{} & \resultcell{} & \resultcell{} & \resultcell{} \\
Prod.-baseline $\times$ anxious & \resultcell{} & \resultcell{} & \resultcell{} & \resultcell{} \\
Prod.-baseline $\times$ positive & \resultcell{} & \resultcell{} & \resultcell{} & \resultcell{} \\
Prod.-integrity $\times$ anxious & \resultcell{} & \resultcell{} & \resultcell{} & \resultcell{} \\
Prod.-integrity $\times$ positive & \resultcell{} & \resultcell{} & \resultcell{} & \resultcell{} \\
Multi-turn interaction mode & \resultcell{} & \resultcell{} & \resultcell{} & \resultcell{} \\
Source complexity & \resultcell{} & \resultcell{} & \resultcell{} & \resultcell{} \\
\bottomrule
\end{tabularx}
\end{table}

\placeholder{If integrity guidance lowers omission but increases disclaimer washing, discuss substitution between failure modes. If its apparent benefit disappears after length matching, report both the total practical effect and the length-adjusted estimate rather than dismissing either.}

\section{RQ4: Persona effects}

Across prompt conditions, the anxious persona produced a \resultcell{} percentage-point change in the primary endpoint relative to neutral (95\% CI \resultcell--\resultcell), while the positive persona produced a \resultcell-point change (95\% CI \resultcell--\resultcell). Persona effects were \resultcell{} across target models and \resultcell{} across system-prompt conditions.

\begin{table}[H]
\centering
\caption{Marginal outcome rates by persona.}
\label{tab:persona-results}
\small
\begin{tabularx}{\textwidth}{@{}Yccc@{}}
\toprule
Outcome & Neutral persona & Anxious risk-averse & Positive risk-seeking/trusting \\
\midrule
High-adverse disclosure (\%) & \resultcell{} & \resultcell{} & \resultcell{} \\
Specificity preservation & \resultcell{} & \resultcell{} & \resultcell{} \\
Minimising language (\%) & \resultcell{} & \resultcell{} & \resultcell{} \\
Benefit amplification (\%) & \resultcell{} & \resultcell{} & \resultcell{} \\
Adverse placement gap & \resultcell{} & \resultcell{} & \resultcell{} \\
Median response tokens & \resultcell{} & \resultcell{} & \resultcell{} \\
Distorted final beliefs (\%) & \resultcell{} & \resultcell{} & \resultcell{} \\
Unsupported/harmful action (\%) & \resultcell{} & \resultcell{} & \resultcell{} \\
\bottomrule
\end{tabularx}
\end{table}

\placeholder{Report whether persona altered substantive disclosure or only style/length. A persona effect should not be interpreted as sycophancy unless it tracks the user's expressed preference in a direction that sacrifices evidence quality. If the anxious condition receives more disclosure but more minimising reassurance, present both rather than assigning one overall label.}

\section{Model heterogeneity}

The evaluated models varied by \resultcell{} percentage points in production-baseline risk-concealment rate. \placeholder{Model name} had the lowest estimated high-adverse omission rate, while \placeholder{model name} had the highest. The integrity prompt reduced the endpoint for \resultcell{} of \resultcell{} models, with model-specific effects ranging from \resultcell{} to \resultcell{} points. Persona sensitivity was strongest for \placeholder{model} and weakest for \placeholder{model}.

\begin{table}[H]
\centering
\caption{Model-specific marginal effects of prompt and persona on the primary endpoint.}
\label{tab:model-effects}
\small
\begin{tabularx}{\textwidth}{@{}Ycccc@{}}
\toprule
Model & Baseline rate & Integrity effect & Anxious effect & Positive effect \\
\midrule
\placeholder{Model A} & \resultcell{} & \resultcell{} & \resultcell{} & \resultcell{} \\
\placeholder{Model B} & \resultcell{} & \resultcell{} & \resultcell{} & \resultcell{} \\
\placeholder{Model C} & \resultcell{} & \resultcell{} & \resultcell{} & \resultcell{} \\
\placeholder{Model D} & \resultcell{} & \resultcell{} & \resultcell{} & \resultcell{} \\
\placeholder{Model E} & \resultcell{} & \resultcell{} & \resultcell{} & \resultcell{} \\
\bottomrule
\end{tabularx}
\end{table}

\placeholder{Avoid ranking models on a single point estimate when intervals overlap or provider versions differ. Include the exact run dates and explain that results are version-specific.}

\section{Single-turn and multi-turn dynamics}

First-turn high-adverse disclosure was \resultcell\% in single-turn scenarios and \resultcell\% in matched multi-turn first responses. Across the complete multi-turn conversation, ever-disclosed rate rose to \resultcell\%, but \resultcell\% of high-adverse facts remained persistently omitted. When the simulator asked directly about the relevant risk category, the target model repaired the omission in \resultcell\% of cases and continued to omit or de-specify it in \resultcell\%.

\begin{table}[H]
\centering
\caption{Multi-turn disclosure and repair outcomes.}
\label{tab:multiturn}
\small
\begin{tabularx}{\textwidth}{@{}Yccc@{}}
\toprule
Outcome & Neutral & Anxious & Positive \\
\midrule
First-turn disclosure (\%) & \resultcell{} & \resultcell{} & \resultcell{} \\
Ever disclosed (\%) & \resultcell{} & \resultcell{} & \resultcell{} \\
Median first-disclosure turn & \resultcell{} & \resultcell{} & \resultcell{} \\
Persistent omission (\%) & \resultcell{} & \resultcell{} & \resultcell{} \\
Repair after direct question (\%) & \resultcell{} & \resultcell{} & \resultcell{} \\
Later minimisation/contradiction (\%) & \resultcell{} & \resultcell{} & \resultcell{} \\
\bottomrule
\end{tabularx}
\end{table}

\placeholder{Distinguish volunteered disclosure from disclosure elicited by the simulator. Show at least one trajectory in which risk is repaired and one in which an initially accurate disclosure is later weakened.}

\section{Belief and action outcomes}

After conversation-only exposure, the user simulator held at least one materially distorted belief in \resultcell\% of multi-turn runs and chose an unsupported or harmful action in \resultcell\%. After receiving the complete fact pool, these rates were \resultcell\% and \resultcell\%, respectively. The within-run action-change rate was \resultcell\%.

High-adverse omission was associated with \resultcell{} higher odds of a distorted belief (95\% CI \resultcell--\resultcell). Misleading framing and de-emphasis showed associations of \resultcell{} and \resultcell. \placeholder{Describe these as simulator-conditioned associations, not estimates of human causal effects.}

\begin{table}[H]
\centering
\caption{Communication outcomes and downstream simulator harm proxies.}
\label{tab:harm}
\small
\begin{tabularx}{\textwidth}{@{}Yccc@{}}
\toprule
State & Distorted belief (\%) & Unsupported action (\%) & Action change after full facts (\%) \\
\midrule
No labelled concealment & \resultcell{} & \resultcell{} & \resultcell{} \\
High-adverse omission & \resultcell{} & \resultcell{} & \resultcell{} \\
Framing only & \resultcell{} & \resultcell{} & \resultcell{} \\
De-emphasis only & \resultcell{} & \resultcell{} & \resultcell{} \\
Multiple failure modes & \resultcell{} & \resultcell{} & \resultcell{} \\
\bottomrule
\end{tabularx}
\end{table}

\section{Robustness analyses}

The direction of the primary integrity contrast was \resultcell{} under fact-order randomisation, \resultcell{} under the length-matched prompt control, \resultcell{} at the alternative output cap, and \resultcell{} when restricted to human labels. The persona result was \resultcell{} across simulator models. Family-leave-one-out estimates ranged from \resultcell{} to \resultcell.

\begin{table}[H]
\centering
\caption{Robustness of the primary prompt and persona effects.}
\label{tab:robustness}
\small
\begin{tabularx}{\textwidth}{@{}YccY@{}}
\toprule
Analysis & Integrity effect & Positive-persona effect & Interpretation \\
\midrule
Preregistered primary & \resultcell{} & \resultcell{} & \placeholder{reference estimate} \\
Fact-order randomised & \resultcell{} & \resultcell{} & \placeholder{position sensitivity} \\
Length-adjusted & \resultcell{} & \resultcell{} & \placeholder{direct versus total effect} \\
Human-labelled subset & \resultcell{} & \resultcell{} & \placeholder{annotation sensitivity} \\
Alternative judge & \resultcell{} & \resultcell{} & \placeholder{judge dependence} \\
Second simulator & \resultcell{} & \resultcell{} & \placeholder{simulator dependence} \\
High-adverse facts only & \resultcell{} & \resultcell{} & \placeholder{materiality sensitivity} \\
\bottomrule
\end{tabularx}
\end{table}

\section{Qualitative error analysis}

Table~\ref{tab:qualitative} should include short, permission-compatible excerpts selected using a pre-declared sampling rule, not only the most dramatic cases. Each example should show the source fact, response transformation, automated label, human judgment, and why the example matters.

\begin{table}[H]
\centering
\caption{Template for qualitative cases.}
\label{tab:qualitative}
\small
\begin{tabularx}{\textwidth}{@{}p{0.13\textwidth}YYYY@{}}
\toprule
Case & Material source fact & Response behaviour & Label & Interpretation \\
\midrule
Q1 & \placeholder{exact adverse fact} & \placeholder{short excerpt/paraphrase} & \placeholder{omission} & \placeholder{selection without falsehood} \\
Q2 & \placeholder{exact quantity/condition} & \placeholder{vague substitution} & \placeholder{specificity loss} & \placeholder{fact present, risk impression weakened} \\
Q3 & \placeholder{adverse and favourable facts} & \placeholder{ordering/allocation pattern} & \placeholder{de-emphasis} & \placeholder{salience effect} \\
Q4 & \placeholder{complete risk evidence} & \placeholder{balanced response} & \placeholder{true negative} & \placeholder{what successful communication looks like} \\
Q5 & \placeholder{ambiguous fact} & \placeholder{contested response} & \placeholder{judge disagreement} & \placeholder{boundary of construct} \\
\bottomrule
\end{tabularx}
\end{table}

\section{Optional detector results}

\begin{placeholderblock}
Delete this section if the detector extension is not completed. If retained, report family-held-out and cross-condition results, calibration, and the incremental value of source and activation access. Do not compare detectors on random response-level splits.
\end{placeholderblock}

\begin{table}[H]
\centering
\caption{Detection performance on held-out scenario families.}
\label{tab:detectors}
\small
\begin{tabularx}{\textwidth}{@{}Ycccc@{}}
\toprule
Detector & Inputs & AUROC & AUPRC & Calibration error \\
\midrule
Response-only black box & Response & \resultcell{} & \resultcell{} & \resultcell{} \\
Source-aware black box & Source + response & \resultcell{} & \resultcell{} & \resultcell{} \\
White-box probe & Activations + response & \resultcell{} & \resultcell{} & \resultcell{} \\
Combined monitor & Source + activations + response & \resultcell{} & \resultcell{} & \resultcell{} \\
\bottomrule
\end{tabularx}
\end{table}

The black-to-white performance boost was \resultcell{} AUROC and \resultcell{} AUPRC. A detector trained on \placeholder{source condition/failure mode} transferred to \placeholder{held-out condition} with \resultcell{} performance, compared with \resultcell{} in-domain. \placeholder{Discuss whether failures favour targeted rather than universal detectors.}
